\begin{landscape}
\small
\renewcommand{\arraystretch}{0.3}

\begin{longtable}{|p{0.3\linewidth}|p{0.4\linewidth}|p{0.2\linewidth}|}
\caption{Session 5 canonical functions: what to implement, what stays the same from Session 4, and where positional, keyword, and mixed calls are introduced.} \\
\hline
\textbf{Section / Function} & \textbf{Code to put inside the function} & \textbf{Session 5 note / modification} \\
\hline
\endfirsthead

\multicolumn{3}{c}%
{{\bfseries \tablename\ \thetable{} -- continued from previous page}} \\
\hline
\textbf{Section / Function} & \textbf{Code to put inside the function} & \textbf{Session 5 note / modification} \\
\hline
\endhead

\hline \multicolumn{3}{|r|}{{Continued on next page}} \\ \hline
\endfoot

\hline \hline
\endlastfoot

\texttt{make\_print\_status(status\_text)} & \begin{minted}[fontsize=\small, linenos, breaklines, xleftmargin=2em]{python}
print(f"[STATUS] {status_text}")
\end{minted}
& Same as Session 4. \\
\hline
\texttt{setup\_application\_list()} & \begin{minted}[fontsize=\small, linenos, breaklines, xleftmargin=2em]{python}
return [
    {"id": "flower1", "petal_length": 1.4, "species": "setosa"},
    ...
]
\end{minted}
& Same role as Session 4 and Session 6: provide the dataset for the lesson. \\
\hline
\texttt{compute\_threshold\_prediction(sample, threshold)} & \begin{minted}[fontsize=\small, linenos, breaklines, xleftmargin=2em]{python}
if sample["petal_length"] < threshold:
    return POSITIVE_LABEL
return NEGATIVE_LABEL
\end{minted}
& Same prediction rule as Session 4 and Session 6. \\
\hline
\texttt{derive\_true\_label(sample)} & \begin{minted}[fontsize=\small, linenos, breaklines, xleftmargin=2em]{python}
if sample["species"] == POSITIVE_LABEL:
    return POSITIVE_LABEL
return NEGATIVE_LABEL
\end{minted}
& Same idea as Session 4 and Session 6. \\
\hline
\texttt{update\_result\_counts(correct, wrong, total, y\_pred\_list, y\_pred, y\_true)} & \begin{minted}[fontsize=\small, linenos, breaklines, xleftmargin=2em]{python}
if y_pred == y_true:
    correct += 1
else:
    wrong += 1
total += 1
y_pred_list.append(y_pred)
return correct, wrong, total, y_pred_list
\end{minted}
& Same counter logic as Session 4 and Session 6. \\
\hline
\texttt{run\_prediction\_loop(dataset, threshold=2.0, print\_each=True)} & \begin{minted}[fontsize=\small, linenos, breaklines, xleftmargin=2em]{python}
for sample in dataset:
    y_pred = compute_threshold_prediction(sample, threshold)
    y_true = derive_true_label(sample)
    correct, wrong, total, y_pred_list = update_result_counts(
        correct, wrong, total, y_pred_list, y_pred=y_pred, y_true=y_true
    )
    if print_each:
        print(...)
\end{minted}
& Same loop structure as Session 4 and Session 6, but now used to teach mixed positional+keyword calls. \\
\hline
\texttt{calculate\_accuracy(correct=0, total=0)} & \begin{minted}[fontsize=\small, linenos, breaklines, xleftmargin=2em]{python}
if total > 0:
    accuracy = (correct / total) * 100
else:
    accuracy = 0.0
return accuracy
\end{minted}
& \textbf{Refactor in S5:} same accuracy logic, but now with keyword-friendly defaults. \\
\hline
\texttt{run\_classifier\_pipeline(threshold=2.0, print\_each=True)} & \begin{minted}[fontsize=\small, linenos, breaklines, xleftmargin=2em]{python}
dataset = setup_application_list()
correct, wrong, total, y_pred_list = run_prediction_loop(
    dataset, threshold=threshold, print_each=print_each
)
accuracy = calculate_accuracy(correct=correct, total=total)
return correct, wrong, total, y_pred_list, accuracy
\end{minted}
& \textbf{New in S5:} renamed wrapper with a clearer purpose and a mixed-call example. \\
\hline
\texttt{print\_summary(correct, wrong, total, y\_pred\_list, accuracy)} & \begin{minted}[fontsize=\small, linenos, breaklines, xleftmargin=2em]{python}
print("Correct:", correct)
print("Wrong:", wrong)
print("Total:", total)
print("Accuracy (%):", round(accuracy, 2))
print("All predictions:", y_pred_list)
\end{minted}
& Same summary idea as Session 4 and Session 6. \\
\hline
\texttt{main()} & \begin{minted}[fontsize=\small, linenos, breaklines, xleftmargin=2em]{python}
run default call
run positional override call
run keyword override call
run reordered keyword call
\end{minted}
& \textbf{New in S5:} demonstrates the same pipeline in several calling styles. \\
\hline
\end{longtable}
\end{landscape}
